\chapter*{Vorwort}
\addcontentsline{toc}{chapter}{Vorwort}

\begin{merksatz}
\boxtitel{Was dieser Band von Band 17 unterscheidet}
Die Statistik aus Band 17 nahm an, dass die erklärenden Größen gegeben sind und die Störungen unabhängig von ihnen. In den Daten, mit denen Volkswirte arbeiten, ist beides regelmäßig falsch.

\begin{center}
{\small
\renewcommand{\arraystretch}{1.52}
\begin{tabular}{@{}P{2.4cm}P{6.2cm}@{}}
\toprule
\textbf{Preise} & entstehen aus Angebot \kursiv{und} Nachfrage --- sie erklären die Menge und werden von ihr erklärt.\\
\textbf{Verträge} & werden gewählt, nicht zugeteilt --- wer viel zum Arzt geht, wählt einen anderen Tarif.\\
\textbf{Erwartungen} & richten sich nach der Politik, die man untersuchen will.\\
\bottomrule
\end{tabular}}
\end{center}

\begin{center}
\begin{minipage}{0.9\textwidth}
\textbf{Die Daten entstehen durch Entscheidungen derer, die man untersucht. Das ist der Gegenstand dieses Bandes.}
\end{minipage}
\end{center}
\end{merksatz}

\begin{merksatz}
\boxtitel{Wo der Band steht}
\begin{center}
{\small
\renewcommand{\arraystretch}{1.5}
\begin{tabular}{@{}P{2.2cm}P{6.4cm}@{}}
\toprule
\textbf{Band 17} & schätzt und testet, bei gegebenen Merkmalen.\\
\textbf{Band 22} & sagt vorher, ohne nach Ursachen zu fragen.\\
\textbf{Band 23} & klärt, wann ein Zusammenhang ursächlich ist.\\
\textbf{Band 25} & liefert die Prozesse, aus denen Zeitreihen bestehen.\\
\textbf{Hier} & werden alle vier gebraucht --- und zwar gleichzeitig.\\
\bottomrule
\end{tabular}}
\end{center}

\begin{center}
\begin{minipage}{0.88\textwidth}
\kursiv{\textbf{Die Ökonometrie ist kein eigenes Gebiet neben diesen, sondern ihre Verbindung unter erschwerten Bedingungen}: kausale Fragen, Beobachtungsdaten, rückgekoppelte Systeme und Zeitreihen, die selten stationär sind. Sie hat dafür eigene Begriffe entwickelt, die in den Nachbarfächern fehlen --- allen voran die Identifikation.}
\end{minipage}
\end{center}
\end{merksatz}

\begin{merksatz}
\boxtitel{Drei Ergebnisse, die man nicht erwartet}
\begin{center}
{\small
\renewcommand{\arraystretch}{1.55}
\begin{tabular}{@{}P{1.0cm}P{7.4cm}@{}}
\toprule
$1$ & Aus Preis- und Mengendaten allein lässt sich \kursiv{weder} die Nachfrage- \kursiv{noch} die Angebotskurve bestimmen --- gleich wie viele Beobachtungen man hat. \kursiv{(Kapitel 2)}\\
$2$ & Zwei voneinander völlig unabhängige Zeitreihen zeigen in den meisten Fällen einen hochsignifikanten Zusammenhang. \kursiv{(Kapitel 11)}\\
$3$ & Ein Modell, das die Vergangenheit gut beschreibt, kann gerade deshalb für die Beurteilung einer Politikänderung untauglich sein. \kursiv{(Kapitel 13)}\\
\bottomrule
\end{tabular}}
\end{center}

\begin{center}
\begin{minipage}{0.88\textwidth}
\kursiv{Das zweite Ergebnis ist das folgenreichste. \textbf{Bei nichtstationären Reihen liegt der Anteil scheinbar bedeutsamer Zusammenhänge bei über siebzig Prozent} --- nicht bei fünf. Jede Zeitreihenauswertung, die das übersieht, produziert Befunde am laufenden Band.}
\end{minipage}
\end{center}
\end{merksatz}

\begin{merksatz}
\boxtitel{Die sechs Teile}
\begin{center}
{\small
\renewcommand{\arraystretch}{1.5}
\begin{tabular}{@{}P{1.4cm}P{3.6cm}P{3.0cm}@{}}
\toprule
\textbf{Teil} & \textbf{Thema} & \textbf{Kapitel}\\
\midrule
I & Das Problem & $1$--$3$\\
II & Querschnittsdaten & $4$--$6$\\
III & Dieselben über die Zeit & $7$--$9$\\
IV & Zeitreihen & $10$--$12$\\
V & Vorhersage und Politik & $13$--$14$\\
VI & Stand & $15$\\
\bottomrule
\end{tabular}}
\end{center}

\schritt{Der Leitfall}
Wie stark sinkt die Zahl der Arztbesuche, wenn Patienten zuzahlen müssen?

\begin{center}
\begin{minipage}{0.88\textwidth}
\kursiv{\textbf{Die Frage hat alle Eigenschaften, um die es hier geht}: Wer krank ist, wählt einen anderen Tarif; wer zuzahlt, geht seltener, aber vielleicht auch dann nicht, wenn er sollte; und die Zuzahlung wird politisch gesetzt, also nicht zufällig. Zugleich gibt es einen randomisierten Versuch, an dem sich alle Verfahren messen lassen --- das amerikanische Krankenversicherungsexperiment der 1970er Jahre.}
\end{minipage}
\end{center}
\end{merksatz}

\begin{warnung}
\boxtitelr{Drei Warnungen}
\begin{center}
{\small
\renewcommand{\arraystretch}{1.55}
\begin{tabular}{@{}P{2.2cm}P{6.4cm}@{}}
\toprule
\textbf{Erste} & Kein Schätzverfahren rettet eine Frage, die nicht identifiziert ist. Das ist eine Eigenschaft des Modells, nicht der Daten.\\
\textbf{Zweite} & Zeitreihen verleiten zu Befunden. Wer nicht auf Stationarität prüft, findet fast alles.\\
\textbf{Dritte} & Eine geschätzte Beziehung gilt für die Politik, unter der sie beobachtet wurde --- nicht notwendig für eine andere.\\
\bottomrule
\end{tabular}}
\end{center}

\begin{center}
\begin{minipage}{0.88\textwidth}
\kursiv{Die erste Warnung ist der eigentliche Beitrag der Ökonometrie zu den Nachbarfächern. \textbf{Identifikation fragt, ob eine Größe aus unbegrenzt vielen Daten überhaupt bestimmbar wäre} --- vor jeder Frage nach Schätzung, Streuung und Test. Band 23 stellte dieselbe Frage mit Graphen; hier wird sie mit Gleichungen gestellt.}
\end{minipage}
\end{center}
\end{warnung}

\begin{merksatz}
\boxtitel{Was dieser Band voraussetzt}
\begin{center}
{\small
\renewcommand{\arraystretch}{1.48}
\begin{tabular}{@{}P{3.0cm}P{1.6cm}P{3.0cm}@{}}
\toprule
\textbf{Werkzeug} & \textbf{Band} & \textbf{gebraucht in}\\
\midrule
Regression, Tests & $17$ & durchgehend\\
Matrizen & $2$ & Kapitel 2--4\\
Kausale Begriffe & $23$ & Kapitel 1--3, 9\\
Prozesse & $25$ & Kapitel 10--12\\
Lernverfahren & $22$ & Kapitel 14\\
\bottomrule
\end{tabular}}
\end{center}
\end{merksatz}
